% UONTEST.TEX - a sample file for the UONTHESIS document style.
% revised by Colin Johnson (Colin.Johnson@nottingham.ac.uk), 2020-02-18
%
\documentclass[11pt]{uonthesis}
\usepackage[style=ieee,backend=biber]{biblatex}
\usepackage[titletoc]{appendix}
\usepackage{graphicx,indentfirst} % put additional packages here
\addbibresource{uontest.bib}
%\bibliographystyle{ieee}   % bibstyle - any commonly used style is acceptable

\renewcommand{\appendixtocname}{List of appendices}
\setlength{\parindent}{2em}
\setlength{\parskip}{1em}

\title{(temp) Using Word Embedding for Citation Recommendation 
in FinTech Scientific Articles}
\author{Ting-Chun Chen}
\prevdegrees{} % if you want to include this
\university{The University of Nottingham}
\degree{Master of Science in Data Science}
\date{September 2023}

%%%%%%%%%%%%%%%%%%%%%%%%%%%%%%%%%%%%%%%%%%%%%%%%%%%%%%%%%%%%%%%%%%%%%%%%%%%%%%%%%%%%
\begin{document}

\begin{frontmatter}
\maketitle
\tableofcontents

\begin{abstract}
According to the guidelines the abstract should not exceed 300 words.
However, it is unlikely that anyone will be counting when you submit.
If you still wish to do so, then better use Emacs \verb|count-words|
command.
\end{abstract}

\begin{acknowledgements}
It is good to thank here your supervisors and any sponsoring bodies,
as well as any family, friends, cats, dogs etc. that have been
supportive during your time at University.
\end{acknowledgements}

%\begin{dedication}
%If you decide to dedicate your thesis to someone, it will be on a
%separate page.
%\end{dedication}

\end{frontmatter}

%%%%%%%%%%%%%%%%%%%%%%%%%%%%%%%%%%%%%%%%%%%%%%%%%%%%%%%%%%%%%%%%%%%%%%%%%%%%%%%%%%%%
%introduction - 2000 words
\chapter{Introduction}
\section{(About Citation)}

Citations are used to demonstrate research background, existing techniques and pieces of evidence for a statement, playing a critical role in scientific writing.
Authors can properly acknowledge the source of information of a statement and avoid plagiarism with an accurate citation.
It also serves as a verification that the idea of this statement is provided and supported by previous studies.
These days, references and citations also include additional information for search engines and citation recommendation systems to recognise the subfields and similar research of this article.

There has been a noticeable increase in the number of scientific articles being published.
3.8 million scientific articles are being published in 2022, according to the Web of Science database.
It is expected that there would be more scientific articles available on the internet.
It can be harder for academics to absorb all the new perspectives with their exact sources.

When looking for reference papers, keywords are commonly used in most academic databases and search engines.
The authors would then review the results of papers from search engines to evaluate their relevance to our research and reliability.
It could cost a huge amount of time and effort to go through full papers, not to mention that the academic database and search engines might miss some relevant papers.

Some academic databases come with citation recommendation tools, which provide relevant articles that share the same category with or are similar to our research.
Citation recommendation tools would have a bias towards published journals, impact factors, times being referred to, and the availability of the journal if it's open to everyone or subscription-only.

\section{Natural Language Processing}

Natural language processing (NLP) refers to the approach of giving computers the ability to understand the meaning and thoughts of words, sentences and articles just like human beings do, building a way of communication between computer and human language.
However, human languages have several characteristics that make it tricky to process them.
The variability and informality, such as different vocabularies, syntax and phrases in different cultures or individuals representing a similar meaning, bring difficulties to applying a general transformations pipeline for information extraction.
A large amount of noise such as mis-spellings, irrelevant contexts and grammar errors confuse not only human readers but also machines and affect the performance of language processing algorithms, not to mention that people might have different understandings of the same sentence which yield even more inaccuracies.
Some advanced circumstances such as sarcasm, irony and the speaker's intention, whose implied meanings are beyond plain text, are more difficult for algorithms to detect and extract straightforwardly. 
Problems in normal vector data such as missing values and lack of labelling occur in natural language as well.

NLP manages to overcome the challenges of processing human language mentioned above and enables machines to analyse, cluster or classify not only vector instances but also texts. 
Advanced applications of NLP, such as sentiment analysis, text clustering, text summarisation, human language generation and improving automatic translation, provide solutions and assistance to a wide range of tasks. 
It reduces our workloads and time-spending and helps improve the interaction between humans and computers by mimicking human language to express in an easy-understanding way for most o the people not only programmers.



% literature review - 4000 words
\chapter{Literature Review}
\section{Text Embedding}
Text embedding methods can be roughly classified by the level of encoding unit (embed in a word, sentence, or article level), supervised or unsupervised, similarity measurement (such as Euclidean distance or cosine similarity) and other specific technique support such as term-based embedding in particular fields and graph-based embedding.
Some advanced methods such as ELMo\cite{peters2018} and GPT\cite{brown2020} use autoregressive models and BERT\cite{niven2019} uses an autoencoder to do bidirectional context encoding.
A more complex structure of neural networks and a bigger training corpus could derive a better and more robust model but a huge amount of time to train models as well.
Hence, another improvement of embedding methods would be the trade-off between better performance and reasonable training time.

\subsubsection{TF-IDF}
Sample text sample text sample text sample text sample text

\subsubsection{LSA Encoding}
Sample text sample text sample text sample text sample text

\subsubsection{Word2Vec Encoding}
Sample text sample text sample text sample text sample text

\subsubsection{Doc2Vec Encoding}
Sample text sample text sample text sample text sample text

\subsubsection{Universal Sentence Encoder}
Sample text sample text sample text sample text sample text

\subsubsection{GloVe}
Sample text sample text sample text sample text sample text

\subsubsection{Long-Short Term Memory}
Sample text sample text sample text sample text sample text

\subsubsection{BERT}
Sample text sample text sample text sample text sample text

\subsubsection{GPT}
Sample text sample text sample text sample text sample text

\subsubsection{ELMo}
Sample text sample text sample text sample text sample text

\subsubsection{Applications of Text Embedding Techniques}
Text embedding techniques are improved for specific objectives, fields and styles and applied to a wide range of tasks, including various issues of tweets analysis\cite{mottaghinia2021}, visualisation in the biomedical field\cite{oubenali2022} and sentiment analysis on movie reviews\cite{sivakumar2021}.

Khatua et al.\cite{khatua2019} identified crisis-related tweets during the 2014 Ebola and 2016 Zika outbreaks with pre-trained Word2Vec and GloVe models.
They found a better classification performance to have a small domain-specific corpus from tweets and scholarly abstracts from PubMed participated in the model.
They also observed a higher accuracy from a higher dimension of word vector and skip-gram model than CBOW.

Lee et al.\cite{lee2017} utilized SentiWordNet 3.0 to analyse the effect of several negative emotions in hotel reviews.
SentiWordNet 3.0\cite{baccianella2010} provided sentiment analysis as classification tasks and word embedding with a frequency-weighted bag-of-words model and the help of WordNet corpus.
Onan\cite{onan2021} presented a sentiment analysis approach to product reviews from Twitter.
This deep-learning-based method applied TF-IDF weighted GloVe to the CNN-LSTM architecture to do word embedding and outperformed conventional deep-learning methods.


\section{Text Similarity}
A pair of similar words should act similarly in most of the features we extracted, while a good similarity metric can aggregate the difference in every feature into a single value, which is comparable between each pair of words.
Besides Euclidean distance and cosine similarity as semantic similarity measurements mentioned above, WordNet also provides path similarity to evaluate similarities of words with a lexical hierarchical structure.

Sentence similarity measurements are more complex and various compared to word similarity since sentences can be considered as combinations of words.
The most straightforward approach is to aggregate words in the sentence as a representation to compare with other sentences, which can be seen as a baseline measurement of sentence similarity.
The steps are to take averages of every word in each sentence, yield a single vector for each sentence and calculate the Euclidean distance or cosine similarity between two sentences with average vectors.
This approach, however, doesn't consider the order of words, while it can have a huge effect on a sentence's meaning when the word order differs.

Several algorithms are proposed to map a sentence into a vector and calculate the similarity with the derived value directly from the embedding process.
Much of them are extended from an existing word embedding method to apply to sentences or even documents, such as Word2Vec, Sent2Vec and Doc2Vec.
Some Approaches consider different levels of embedding at the same time, such as word embedding and sentence embedding of BERT.

\subsubsection{Lexical Relation and WordNet}
Sample text sample text sample text sample text sample text


\section{Text Embedding for Search Engine}
With the increase of documents and web pages, traditional keywords-based search engines are thought to be powerless to correctly look for users' requirements.
Text embedding techniques are applied to search engines to provide machine-readable web pages and semantic annotations to the algorithm to yield a more accurate search result.
Many websites are applying text embedding to their search engines, including Google Search, Microsoft's Bing Search, Amazon Search and many e-commerce websites and platforms.
Not to confuse the search engine using text embedding and the search engine using semantic web search language, in this paper, we refer to the semantic search as the searching approach with text embedding or other semantic retrieval techniques. 

Regarding the search engines for scientific articles, Eisenberg et al.\cite{eisenberg2017} proposed a semantic search for Biogeochemical literature that undergoes paper research by comparing the concepts extracted from queries and academic literature.
However, the concept extraction component in the workflow of this research is annotated by domain experts, which can be done automatically by NLP approaches mentioned in the future directions chapter.
Fang et al.\cite{fang2018} performed a biomedical article-searching approach called Semantic Sequential Dependence Model (SSDM) that combined semantic information retrieval techniques and the traditional SDM model.
Word embedding techniques of the Neural Network Language Model(NNLM)\cite{bengio2000}, Log-Bilinear Language Model(LBL)\cite{mnih2007} and Word2vec are applied to the literature corpus to find synonyms by KNN classification algorithm and generate a domain-specific thesaurus.
The thesaurus is then utilized to extend query keywords and the SDM played an important role in the combination strategy of the extended keywords for further comparison to documents.

Compare to the keyword embedding approaches, applying sentence embedding or other wider-level embedding techniques to search engines can better preserve information for query but be more challenging at the same time.
Palangi et al.\cite{palangi2016} managed to embed semantic information at a sentence level by using LSTM-RNN (Long Short-Term Memory - Recurrent Neural Network) framework to do web document retrieval, and compare the summarisation sentence vector and query sentence vector to yield the best searching result.


\section{Searching-Based Citation Recommendation}

Some academic databases come with citation recommendation tools, which provide relevant articles that share the same category with or are similar to our research.
Citation recommendation tools would yield different results resulting from not only different algorithms they used but also candidate papers' published journals, impact factors, times being referred to, and the availability if it's open to everyone or subscription-only.

Citation recommendation methods are always built with three main stages: (1) Generate candidate citations from the publication database (2) Create a recommendation list by ranking the candidate citations (3) Evaluate the accuracy of our recommendation system.\cite{ma2020}
Text embedding techniques can be included generally in step 1: generation of candidate citation, that is, applying text embedding to filter out candidate citations from the publication database that better relate to keywords provided by users or meet users' needs.

Since the diversity of terms and patterns between scientific articles in different domains, researchers would tend to train a specific model with their domain data or use transfer learning to fine-tune the model parameters.
Tshitoyan et al.\cite{tshitoyan2019} used modified Word2Vec embedding to successfully capture complex materials science concepts, without any additional chemistry knowledge insertion, and extract knowledge and relationship from scientific literature.
Zhang et al.\cite{zhang2022} used 15 text representation models, including 6 term-based methods, 2 word embedding methods, 3 sentence embedding methods, 2 document embedding methods and 2 BERT-based methods, to construct an article recommendation system in the biomedical field.
They found BERT and BioSenVec, a Sent2Vec model trained on PubMed corpus, outperformed most of the online and offline citation recommendation systems and an improvement in BERT-based methods after fine-tuning to learn users' preferences.

Wang et al.\cite{wang2022} proposed a sentence-level citation recommendation system called SenCite that used CNN to recognise candidate citation sentences and FastText as the word embedding method to extract information from texts, without summarising an article into sentences.
They evaluated the performance of SenCite with several evaluation metrics such as modified reciprocal rank, average precision and normalized discounted cumulative gain and human experts verification.
The SenCite is shown to outperform most of the embedding methods and yield a comparable accuracy to BERT.
The human experts also stated that SenCite provided the best top-1 citation recommendation.

% material and method - 3000 words
\chapter{Material and Method}
\section{Data Collection}
We use the Web of Science[-citation-] search engine and database to find the fintech articles and their journals.
Reference papers were collected by the WOS search engine with rules listed as follows:
\begin{itemize}
    \item Topic: "\textbf{Fintech}" or "\textbf{Financial technology}" or "\textbf{digital finance}" or "\textbf{electronic banking}" or "\textbf{cryptocurrency}" or "\textbf{Blockchain}"
    \item Document Types: \textbf{Article} or \textbf{Proceeding Paper} or \textbf{Review Article}
    \item Languages: English
\end{itemize}

We got 66,974 results on 23rd July 2023 from the Web of Science Core Collection and the top 10 publications with the most article are as follow:
\begin{table}[]
\caption{Table 1. The top 10 publications with the most fintech articles.}
\centering
    \begin{tabular}{|lc|lc|}
    \hline
    Journal                         & \# of papers & Journal                                      & \# of papers \\ \hline
    IEEE ACCESS                     & 1342         & APPLIED SCIENCES-BASEL                       & 404          \\
    SUSTAINABILITY                  & 1207         & ENVIRONMENTAL SCIENCE AND POLLUTION RESEARCH & 377          \\
    SENSORS                         & 505          & ENERGIES                                     & 374          \\
    IEEE INTERNET OF THINGS JOURNAL & 489          & ELECTRONICS                                  & 344          \\
    JOURNAL OF CLEANER PRODUCTION   & 415          & TECHNOLOGICAL FORECASTING AND SOCIAL CHANGE  & 326          \\ \hline
    \end{tabular}
    \end{table}

Due to the time restriction mentioned in the following section, we utilise the fintech articles from a single Journal \textit{IEEE Access} as the raw data of this research.
The fintech articles are exported to \texttt{EndNote 20}\footnote{https://endnote.com/} and the \texttt{Find Full Text} feature is used to look for their full text by their Accession numbers including Digital Object Identifier (DOI), WOS ID and PubMed ID.
This feature enables us to automatically batch-download loads of reference papers as PDF files if the journal is available in Open Access resources.

\section{Data Structuring}
We convert the unstructured PDF file into a structured database consisting of pairs of a citation sentence and its full-text reference paper.
Firstly, we extract the main content with citations and the bibliography of the articles from PDF files to a plain text file and a dictionary, respectfully.
This process is to remove non-organised and less-informed texts, such as page numbers, article titles, journal names, images, math equations and paragraph/section labels, to make sure the sentences in the main content can be fully extracted for further analysis.
The bibliography of the articles is also extracted and the information of each reference paper is well-organised to a dictionary.
The reference papers extracted from an article are serialised by \texttt{pickle}\footnote{https://docs.python.org/3/library/pickle.html} module to a list of dictionaries.
The citation sentence can then look up to its corresponding reference paper via the information in this reference dictionary and be matched with the full text.
The Python tool \texttt{PyMuPDF}\footnote{https://pypi.org/project/PyMuPDF/} is used to obtain the style and location of lines, spans, and sections in a PDF page.
The patterns including font size, style and colour are set manually and utilised to recognise if the texts in the current span are the main content, bibliography, or non of both.
Hence, the design and layout of a journal would highly affect the criteria of the extraction.
We would have to customise the process of data extraction from one journal to another, which takes a lot of time and effort and we're not yet able to program the automatic recognition.

Secondly, The sentences that are inspired by reference papers, called citation sentences, are extracted and organised with the reference number.
We split the full texts into sentences and recognise the citation sentences by square brackets with a number in them.
The citation sentence and reference number would be organised and serialised into a list.
Last but not least, the citation sentence and the reference paper are matched together to form the database for this research.
The reference information of the citation sentence can be found via the reference number and further match the citation sentence to the plain text of the reference paper.
The citation sentences from all of the articles and their reference texts are saved as a list of dictionaries.
The whole dataset is used for unsupervised models, while a 7:3 ratio is applied to split the database into the training set and the test set for supervised models.

\section{Data Preprocessing}
Data preprocessing includes every approach to convert the database to the required input of every embedding model.
Among all the preprocessing steps, \texttt{unidecode} is priorly applied to texts to convert special characters to the closest English characters, such as \textit{\o}, \textit{\^o}, \textit{\~o}, \textit{\=o}, \textit{\'o}, \textit{\u o}, \textit{\v o}, \textit{\oo}, \textit{\oe} and ligature such as \textit{fi}, \textit{fl} and \textit{ff}
Sentence tokenisation is applied to separate the whole article into sentences basically by the period mark.
Word tokenisation is applied to sentences that split them into words roughly by space.
Stop words are removed from the word tokens and lemmatisation is applied to the word tokens to reduce inflected forms of a word into the word's lemma to treat them as a single word since they contain a similar meaning even if they look different.
The above process is implemented by Natural Language Toolkit, \texttt{nltk}.

\section{Experiments}


% results - 3500 words

% conclusion and discussion - 3000 words



% \begin{table}
% \caption{Tables usually have caption above.}
% \centering
% \begin{tabular}{c|cc}\hline\hline
% Put & a table & here \\
% \hline
% 1 & 2 & 3 \\
% 1 & 2 & 3 \\
% 1 & 2 & 3 \\
% 1 & 2 & 3
% \end{tabular}
% \end{table}


\subsection{Second Level Heading}

Sample text sample text sample text sample text sample text sample
text sample text sample text sample text sample text sample text
sample text sample text sample text sample text sample text.

% \begin{figure}
% \centering
% \includegraphics[width=40mm]{kermit.jpg}
% \caption{You can use {\tt includegraphics} command with the {\tt
%     graphicx} style option to add images.}
% \end{figure}


\subsubsection{Third Level Headings Do Not Appear in the Table of Contents}

Sample text sample text sample text sample text sample text sample
text sample text sample text sample text sample text sample text
sample text sample text sample text sample text sample text~\cite{Lam94}

%%%%%%%%%%%%%%%%%%%%%%%%%%%%%%%%%%%%%%%%%%%%%%%%%%%%%%%%%%%%%%%%%%%%%%%%%%%%%%%%%%%%
%\bibliography{uontest} % filename (before .bib) of your BibTeX database
\printbibliography

%%%%%%%%%%%%%%%%%%%%%%%%%%%%%%%%%%%%%%%%%%%%%%%%%%%%%%%%%%%%%%%%%%%%%%%%%%%%%%%%%%%%
\begin{appendices}

\chapter{Supplementary Materials I}

Sample text sample text sample text sample text sample text sample
text sample text sample text sample text sample text sample text
sample text sample text sample text sample text sample text.

\chapter{Supplementary Materials II}

Sample text sample text sample text sample text sample text sample
text sample text sample text sample text sample text sample text
sample text sample text sample text sample text sample text.

\end{appendices}


\end{document}
